%% ============================================================
%%  Chapter 1 – Introduction
%% ============================================================
\chapter{Introduction}
\label{ch:intro}

\begin{infobox}[Chapter Overview]
This chapter motivates the need for explainable, privacy-compliant Medical
Visual Question Answering (VQA), states the research questions addressed by
\medxplain, and summarises the core contributions of this thesis.
\end{infobox}

%% ─────────────────────────────────────────────────────────────
\section{Background and Motivation}
\label{sec:motivation}

Medical imaging occupies a central role in modern clinical decision-making.
Chest radiography alone accounts for more than 2 billion examinations
worldwide each year~\citep{who2020xray}, with demand accelerating as
populations age and portable imaging technology proliferates. Yet the supply
of trained radiologists grows far more slowly. Studies consistently report
that radiologist workload has doubled over the past decade, contributing to
fatigue, longer reporting times, and measurably higher diagnostic error
rates~\citep{reiner2012radiology}.

Artificial intelligence offers a compelling complement to human expertise.
Deep learning models have achieved radiologist-level performance on isolated
classification tasks such as detecting pneumonia~\citep{rajpurkar2017chexnet}
or screening for diabetic retinopathy~\citep{gulshan2016retina}. However,
two fundamental gaps prevent these systems from entering routine clinical use:

\begin{enumerate}[label=\arabic*.]
  \item \textbf{Lack of explainability.} Clinicians and regulators require
        insight into \emph{why} a model reaches its conclusion.
        The European Union's AI Act (Article~13) and the US FDA's action plan
        for AI/ML-based Software as a Medical Device both mandate transparency
        as a prerequisite for high-risk AI deployment~\citep{euaiact,fda2021ai}.
  \item \textbf{Incompatibility with hospital data infrastructure.}
        Hospital imaging data is stored in the \dicom{} format, legally
        protected under HIPAA~\citep{hipaa} and GDPR~\citep{gdpr}, and must
        never leave the institution's network. Existing research systems are
        trained on curated public datasets and offer no pathway for ingesting
        private, PHI-bearing DICOM files.
\end{enumerate}

\medxplain{} addresses both gaps simultaneously by combining a principled
privacy-engineering layer with integrated explainability and a clinical-grade
user interface.

%% ─────────────────────────────────────────────────────────────
\section{Problem Statement}
\label{sec:problem}

We articulate the core research problem as follows:

\begin{warningbox}[Research Problem]
Given a medical image $\mathbf{x} \in \mathbb{R}^{H \times W \times C}$
sourced from a hospital PACS system and a natural-language clinical
question $q$, produce: (i)~an accurate answer $a$; (ii)~a spatial
explanation $\mathbf{E} \in \mathbb{R}^{H \times W}$ highlighting the image
regions that justify $a$; (iii)~a confidence score $\hat{c} \in [0,1]$;
all while guaranteeing that no Protected Health Information is retained in
any intermediate or final artefact.
\end{warningbox}

This formulation extends the classic VQA setting~\citep{antol2015vqa} with
two hard constraints: \textbf{privacy} (HIPAA/GDPR compliance) and
\textbf{explainability} (regulatory-grade visual attribution).

%% ─────────────────────────────────────────────────────────────
\section{Research Questions}
\label{sec:rqs}

This thesis investigates three research questions:

\begin{description}
  \item[\textbf{RQ1}] Can a multimodal VQA model trained on de-identified
        hospital chest radiographs achieve accuracy exceeding 65\% on
        closed-form clinical questions, as specified in the Cahier des
        Charges?
  \item[\textbf{RQ2}] How can PHI compliance (HIPAA/GDPR) and
        explainability (Grad-CAM) be engineered as \emph{first-class
        architectural concerns} rather than post-hoc additions?
  \item[\textbf{RQ3}] Does a two-phase training strategy (vision-language
        alignment followed by LoRA fine-tuning) outperform single-phase
        end-to-end training under a limited hospital-GPU budget?
\end{description}

%% ─────────────────────────────────────────────────────────────
\section{Proposed Solution: \medxplain}
\label{sec:solution}

\medxplain{} is a modular, hospital-deployable Medical VQA framework with
the following architectural pillars:

\begin{itemize}
  \item \textbf{Privacy-first data pipeline.} Raw DICOM studies are loaded,
        stripped of all 18 HIPAA Safe Harbour identifiers, and converted to
        normalised float32 tensors before any model ever sees the data
        (\S\ref{sec:dicom}--\ref{sec:phi}).

  \item \textbf{Pluggable vision-language architecture.} A BiomedCLIP
        vision encoder~\citep{zhang2023biomedclip} is fused with a BioGPT
        language decoder~\citep{luo2022biogpt} via a trainable
        cross-attention module, all swappable via YAML configuration.

  \item \textbf{Two-phase training.} Phase~1 aligns the visual and
        textual embedding spaces with NT-Xent contrastive
        loss~\citep{chen2020simclr} on frozen backbones. Phase~2
        fine-tunes for closed-form VQA using LoRA adapters~\citep{hu2021lora}
        (rank~=~8) on the language decoder, reducing trainable parameters
        by ${\sim}99\%$.

  \item \textbf{Integrated explainability.} Every prediction is accompanied
        by a Grad-CAM~\citep{selvaraju2017gradcam} heatmap overlay rendered
        in real-time within the Gradio web interface.

  \item \textbf{Air-gapped deployment.} The system operates entirely
        on-premises. \code{allow\_cloud\_upload: false} is enforced in the
        hospital configuration, and all core functionality requires no
        external network access.
\end{itemize}

%% ─────────────────────────────────────────────────────────────
\section{Contributions}
\label{sec:contributions}

The original contributions of this thesis are:

\begin{enumerate}[label=\textbf{C\arabic*.}]
  \item \textbf{Hospital-Ready Data Pipeline (Chapter~\ref{ch:method},
        \S\ref{sec:data_infra}).}
        A production-grade DICOM ingestion pipeline with HIPAA/GDPR-compliant
        pseudonymisation, VOI-LUT pixel normalisation, and a tamper-evident
        audit trail. Implemented in \code{data\_ingestion/dicom\_pipeline.py},
        \code{phi\_scrubber.py}, and \code{report\_parser.py}.

  \item \textbf{Two-Phase Training Protocol
        (Chapter~\ref{ch:method}, \S\ref{sec:training}).}
        NT-Xent contrastive pre-alignment followed by LoRA VQA fine-tuning
        with automatic mixed precision and gradient checkpointing.
        Implemented in \code{training/train\_vqa\_two\_phase.py}.

  \item \textbf{Integrated Explainability
        (Chapter~\ref{ch:method}, \S\ref{sec:explainability}).}
        Grad-CAM explanations embedded in a privacy-safe clinical UI,
        displayed side-by-side with model answers and confidence scores.

  \item \textbf{Clinical Validation Framework
        (Chapter~\ref{ch:results}, \S\ref{sec:clinical_val}).}
        A radiologist-comparison module that computes Cohen's Kappa, flags
        discrepancies, and generates PHI-free PDF reports for manual review.

  \item \textbf{Extended Evaluation Suite
        (Chapter~\ref{ch:results}, \S\ref{sec:quant}).}
        BERTScore F1 and per-pathology Clinical F1 added alongside
        standard BLEU-4 and exact-match accuracy.
\end{enumerate}

%% ─────────────────────────────────────────────────────────────
\section{Thesis Structure}
\label{sec:structure}

The remainder of this thesis is organised as follows.
Chapter~\ref{ch:lit} reviews related work in medical VQA, vision-language
models, parameter-efficient fine-tuning, and medical AI explainability.
Chapter~\ref{ch:method} describes the complete \medxplain{} methodology,
including the data infrastructure, model architecture, and training protocol.
Chapter~\ref{ch:results} presents experimental results and ablation studies.
Chapter~\ref{ch:discussion} discusses strengths, limitations, and ethical
considerations. Chapter~\ref{ch:conclusion} concludes and outlines future
directions. Appendices provide the repository map, configuration reference,
and radiologist validation form.
